\newpage
\section{first-order perturbative approach for two absorbers}\label{sec:osc-two-absorber-o1}

we consider the case of two atoms, located at $x_1=a_1$ and $x_2=a_2$. The perturbation Hamiltonian is the sum of two terms:
\begin{align}
  H_I = H_I^{(1)}(X-x_1) +H_I^{(2)}(X-x_2)
\end{align}
For the initial test the interaction Hamiltonian caa be chosen as
\begin{align}
  H_I(y)=V_0 e^{-\frac{y^2}{2\sigma_V^2}}
\end{align}
the initial state is given by the alpha particle wave packet and both oscillators in their ground state:
\begin{align}
  |\psi_i\rangle = \int\limits_{-\infty}^{\infty}dP C(P)|P\rangle\otimes|0(a_1)\rangle\otimes|0(a_2)\rangle
\end{align}
here $C(P)$ is the projectile momentum coefficient of Part~\ref{part:general-tdpt}, Sect.~\ref{sec:proj-wavepacket}, given by (\ref{eq:proj-CP}) and, for $P_\alpha\gg\hbar/\sigma_\alpha$, approximated by (\ref{eq:proj-CP-approx}).
the final state where  oscillators are excited has the form:
\begin{align}
  |\psi_f\rangle = |P_f, n_1, n_2\rangle = |P_f\rangle\otimes|n_1(a_1)\rangle\otimes|n_2(a_2)\rangle
\end{align}
we calculate the first order transition amplitude to find the probability of one-atom excitation (note that the first order calculation can predict only one atom excitation)

starting point: the general formula (\ref{eq:gen-first-order-amplitude})
\begin{align}
  {\cal A}_{fi}^{(1)} & =\sum_\gamma {c_\gamma^{(i)}} \langle \gamma_f|H_I|\gamma\rangle  {\cal E}(E_{\gamma_f},E_\gamma)\label{eq:osc-first-order-amplitude-restated}
\end{align}

We apply it for our situation
We start from the general calculation
\begin{align}
  \langle \gamma_1|H_I|\gamma_2\rangle = \langle n^{(1)}_1(a_1)n_1^{(2)}(a_2);P_1|H_{I}^{(1)}+H_{I}^{(2)}|n_2^{(1)}(a_1)n_2^{(2)}(a_2);P_2\rangle
\end{align}
the one oscillator interaction matrix element will be denoted as
\begin{importantbox}
  \begin{align}
    \langle n_1(a) P_1|H_I|n_2(a) P_2\rangle={\cal V}(n_1,n_2; P_1, P_2; a)\label{eq:osc-V-notation}
  \end{align}
\end{importantbox}
and they are calculated in Sect \ref{sec:osc-V-matrix-element}
\begin{importantbox}
  \begin{align}
    {\cal V}(n_1,n_2;P_1,P_2;a)&= \langle n_1(a) P_1|H_I|n_2(a) P_2\rangle  \nonumber\\
    &= \frac{V_0}{2\pi\hbar}\sqrt{{2\pi\sigma_V^2}} e^{-\frac{i}{\hbar}a(P_1-P_2)}e^{-\frac{(P_1-P_2)^2\sigma_V^2}{2\hbar^2}}\sqrt{\frac{n_<!}{n_>!}} \left(-\frac{i (P_1-P_2)  }{\sqrt{2}\hbar\alpha}\right)^{n_> - n_<} e^{-\frac{(P_1-P_2)^2 }{4\hbar^2\alpha^2}} \, L_{n_<}^{(n_> - n_<)}\!\left(\frac{(P_1-P_2)^2 }{2\hbar^2\alpha^2}\right)\label{eq:osc-V-n1n2-restated}
  \end{align}
  where $n_<$ and $n_>$ are smaller/larger between $n_1$ and $n_2$
  For $P_1=P_2$, the expression containing the power of $P_1-P_2$
  should be evaluated through its limiting value.  Using the
  orthonormality of the harmonic-oscillator eigenfunctions gives
  \begin{align}
    {\cal V}(n_1,n_2;P,P;a)
    = \frac{V_0}{2\pi\hbar}\sqrt{2\pi\sigma_V^2}\,
    \delta_{n_1n_2}.\label{eq:osc-V-n1n2-diag-restated}
  \end{align}
\end{importantbox}
with the results above the matrix element $ \langle \gamma_1|H_I{(i)}|\gamma_2\rangle$ can be written as
\begin{align}
  \langle \gamma_1|H_I|\gamma_2\rangle = \delta_{n_1^{(2)}n_2^{(2)}}{\cal V}(n_1^{(1)},n_2^{(1)};P_1,P_2,a_1)+\delta_{n_1^{(1)}n_2^{(1)}}{\cal V}(n_1^{(2)},n_2^{(2)};P_1,P_2,a_2)
\end{align}
For the calculation of the transition amplitude we have
\begin{align}
  &\gamma\rightarrow |P;0(a_1),0(a_2)\rangle,\qquad E_\gamma = \frac{P^2}{2M}+\hbar\omega\\
  & \gamma_f\rightarrow|P_f;n_f^{(1)},n_f^{(2)}\rangle\qquad E_{\gamma_f} = \frac{P_f^2}{2M}+\hbar\omega(n_f^{(1)}+n_f^{(2)}+1)
\end{align}
and the transition amplitude becomes
\begin{importantbox}
  \begin{align}
    {\cal A}_{fi}^{(1)} (P_f;n_f^{(1)},n_f^{(2)})= \int\limits_{-\infty}^{\infty}dPC(P)\left[\delta_{n_f^{(2)}0}{\cal V}(n_f^{(1)},0;P_f,P;a_1)+\delta_{n_f^{(1)}0}{\cal V}(n_f^{(2)},0; P_f,P;a_2)\right] {\cal E}(E_{\gamma_f},E_\gamma)
  \end{align}
\end{importantbox}
the factor ${\cal E}$ is calculated in Sect. \ref{sec:app-stable-first-order} and is
\begin{importantbox}
  \begin{align} {\cal E}(E_{\gamma_f},E_\gamma)=\frac{e^{-i\tau E_{\gamma_f}}-e^{-i\tau E_\gamma}}{E_{\gamma_f}-E_\gamma} = -i\tau
    e^{-i\tau(E_{\gamma_f}+E_\gamma)/2}        \operatorname{sinc} \left[
      \frac{\tau(E_{\gamma_f}-E_{\gamma})}{2}
    \right].\qquad \tau = \frac{t_f-t_i}{\hbar}\label{eq:osc-E-factor-restated}
  \end{align}
\end{importantbox}

The matrix element $ {\cal V}$ has a simplified form since one of the indices is zero
\begin{importantbox}
  \begin{align}
    {\cal V}(n_f,0;P_f,P;a)&= \langle n_f(a) P_f|H_I|0 P\rangle  \nonumber\\
    &= \frac{V_0}{2\pi\hbar}\sqrt{{2\pi\sigma_V^2}} e^{-\frac{i}{\hbar}a(P_f-P)}e^{-\frac{(P_f-P)^2\sigma_V^2}{2\hbar^2}}\sqrt{\frac{1}{n_f!}} \left(-\frac{i (P_f-P)  }{\sqrt{2}\hbar\alpha}\right)^{n_f} e^{-\frac{(P_f-P)^2 }{4\hbar^2\alpha^2}} \label{eq:osc-V-n0}
  \end{align}
  For $P_1=P_2$, the expression containing the power of $P_f-P$
  should be evaluated through its limiting value.  Using the
  orthonormality of the harmonic-oscillator eigenfunctions gives
  \begin{align}
    {\cal V}(n_f,0;P,P;a)
    = \frac{V_0}{2\pi\hbar}\sqrt{2\pi\sigma_V^2}\,
    \delta_{n_f0}.\label{eq:osc-V-n0-diag}
  \end{align}
  and $\alpha$ is
  \begin{align}
    \alpha=\sqrt{\frac{m\omega}\hbar}
  \end{align}
\end{importantbox}
The numerical observable is the differential density probability
\begin{importantbox}
  \begin{align}
    {\cal P}(P_f, n_f^{(1)}, n_f^{(2)}) = \left|{\cal A}_{fi}^{(1)} (P_f;n_f^{(1)},n_f^{(2)})\right|^2
  \end{align}
\end{importantbox}

\begin{importantbox}
  in the numerical calculation there are two regimes:
  \begin{enumerate}
    \item $M\gg m \sim 1$ au. The velocity $v$ of the alpha particle is of the order of 10 au (it has to be smaller that $c = 137.036$ au). In that case the momentum $P_\alpha$ is very large; there will be two regions of interest for $P_f$, in two disconnected stripes around $\pm P_\alpha$, with width of the order of $\sqrt{\frac{\hbar^2 }{\sigma_V^2 + \frac1{2\alpha^2}}}$

    \item $M\sim m\sim1$ au. In that case the velocity of the alpha particle $v$ remains the same, but the momentum will be much lower, and the two stripes from the previous point become a single stripe.
  \end{enumerate}
\end{importantbox}